\chapter{Pelvic Organ Prolapse}

\section*{Definition}
Descent of one or more pelvic organs (uterus, bladder, rectum, bowel, or vaginal vault) into or beyond the vaginal canal due to weakened pelvic floor support structures.

\section*{What Happens in Your Body}
Pelvic organ prolapse results from compromise of the pelvic floor musculature (levator ani complex) and connective tissue (endopelvic fascia, uterosacral ligaments). Oestrogen maintains collagen and elastin in pelvic support tissues; declining oestrogen after menopause reduces tissue strength and elasticity. Vaginal childbirth (particularly multiple or traumatic deliveries) stretches and damages the levator ani and pudendal nerve. Increased intra-abdominal pressure (obesity, chronic cough, constipation, heavy lifting) accelerates prolapse progression. Prolapse is graded I-IV based on descent relative to the hymen.

\section*{Causes}
\begin{itemize}
  \item Menopause (oestrogen deficiency weakening connective tissue)
  \item Vaginal childbirth (especially multiple, instrumental, or macrosomic deliveries)
  \item Advanced age
  \item Obesity (increased intra-abdominal pressure)
  \item Chronic constipation or straining
  \item Chronic cough (smoking, COPD, asthma)
  \item Heavy lifting (occupational or recreational)
  \item Connective tissue disorders (Ehlers-Danlos, Marfan)
  \item Prior pelvic surgery (hysterectomy)
  \item Genetic predisposition (family history of prolapse)
\end{itemize}

\section*{When to See a Doctor}
See your doctor if:
\begin{itemize}
  \item Symptoms are severe or getting worse over time
  \item Sensation of a bulge or pressure in the vagina, pelvic heaviness, urinary or faecal incontinence, or difficulty emptying bladder/bowel
  \item You have other concerning symptoms such as fever, weight loss, or bleeding
\end{itemize}

\section*{Related Symptoms}
\begin{itemize}
  \item Pelvic pain
  \item Urinary incontinence
  \item Urinary frequency
  \item Back pain
  \item Painful sex
\end{itemize}

\textbf{Important:} If this symptom persists, your doctor will check for serious underlying causes.

\section*{Self-Care / Home Management}
\begin{itemize}
  \item Rest and avoid activities that make it worse.
  \item Maintain adequate hydration and a balanced diet.
  \item Over-the-counter remedies may provide symptomatic relief (consult a pharmacist or doctor).
  \item Monitor symptoms and seek professional advice if they persist or worsen.
  \item Avoid self-medicating with prescription medications without medical guidance.
\end{itemize}

\section*{How Your Doctor Will Evaluate You}
\begin{itemize}
  \item A thorough history and physical examination are the first steps in evaluation.
  \item Laboratory tests (blood work, urinalysis) may be ordered based on clinical suspicion.
  \item Imaging studies (X-ray, ultrasound, CT, MRI) may be indicated when structural pathology is suspected.
  \item Specialist referral is arranged when the diagnosis remains uncertain or requires specific expertise.
\end{itemize}

\section*{Treatment Options}
\begin{itemize}
  \item Treatment targets the underlying cause identified through diagnostic evaluation.
  \item Supportive care and symptom management are provided while awaiting definitive therapy.
  \item Medications, lifestyle modifications, and physical therapies may be prescribed as appropriate.
  \item Follow-up ensures treatment effectiveness and allows timely adjustments to the care plan.
\end{itemize}

\clearpage